% ============================================
% Johns Hopkins Statement Template
% Inclusive Learning Environment – Christopher Newport University
% ============================================

\documentclass[11pt,letterpaper]{article}

\usepackage{geometry}
\usepackage{microtype}
\usepackage[hidelinks]{hyperref}
\usepackage{setspace}
\geometry{left=25mm,right=25mm,top=25mm,bottom=25mm}
\setstretch{1.05}

\begin{document}

\begin{center}
    {\Large \textbf{Statement on Fostering an Inclusive Learning Environment}}\\[6pt]
    \textit{Decory Edwards}
\end{center}

My approach to teaching economics begins with the conviction that every student, regardless of background, prior preparation, or lived experience, should feel welcomed, respected, and capable of mastering quantitative and conceptual tools. In a liberal-arts setting like Christopher Newport University, inclusion means designing courses where rigor and accessibility reinforce each other, and where students see economic analysis as relevant to their own experiences and aspirations.

I incorporate diverse perspectives into course materials and discussions. In courses such as \textit{Economics of Inequality and Tax Policy} or \textit{Macroeconomic Theory}, students examine data that reflect gender, racial, and regional variation in economic outcomes and explore how institutional differences shape opportunity. By linking empirical results to historical and social context, I help students recognize economics as a discipline that both explains and challenges inequality. Students are encouraged to connect theory with their own experiences—for example, analyzing saving behavior in different household settings or exploring policy trade-offs faced by small communities.

Mentorship and research inclusion are central to my teaching philosophy. I invite undergraduates to join research projects, providing structured guidance on data management, replication, and presentation. For students underrepresented in economics, I offer individualized coaching on research design and graduate study preparation. I also collaborate with student groups and departmental initiatives to host speakers from diverse backgrounds and organize workshops that demonstrate how economic reasoning can inform civic engagement and leadership.

I view inclusive teaching as an iterative practice. After each semester, I review participation data and anonymous student reflections to identify patterns of engagement and to refine assignments or pacing. I also seek peer observation and exchange strategies with colleagues to strengthen our collective approach to inclusive pedagogy.

Ultimately, fostering inclusion means creating an environment where all students can connect rigorous economic thinking with empathy, curiosity, and a sense of belonging. My goal is for students to leave my classroom confident not only in their analytical skills, but in their voice as contributors to a broader conversation about policy, equity, and prosperity.

\end{document}
